\section{서론}

현대의 온라인 디지털 서비스는 사용자의 행동, 위치, 상호작용, 상태 변화, 사건 발생을 시간 순서대로 기록하는 대규모 telemetry 데이터를 생성한다. 이러한 데이터는 단순한 정적 로그나 최종 결과표가 아니라, 시간에 따라 이동하는 객체, 변화하는 공간적 조건, 그리고 특정 시점에 발생하는 사건이 결합된 시공간 데이터의 성격을 가진다. 온라인 멀티플레이 게임은 이러한 telemetry 데이터가 특히 밀도 높게 생성되는 대표적인 환경이다. 수많은 플레이어가 같은 가상 공간 안에서 동시에 이동하고, 자원을 탐색하며, 다른 플레이어와 상호작용하고, 경기 규칙에 의해 변화하는 환경 조건에 대응하기 때문이다.

게임 산업은 2025년 기준 약 1,888억 달러 규모와 약 35.8억 명 수준의 플레이어 기반을 가진 대규모 디지털 환경으로 성장하였다~\cite{newzoo2025globalgames}. 이 가운데 배틀로얄 장르는 대규모 온라인 경쟁, 실시간 의사결정, 관전자 중심 콘텐츠, e-sports와 streaming 문화의 확산과 함께 중요한 장르적 위치를 차지해 왔다. 배틀로얄 게임은 다수의 플레이어 또는 팀이 하나의 넓은 전장에 동시에 진입한 뒤, 제한된 자원을 확보하고, 다른 플레이어와 경쟁하면서 마지막까지 생존하는 것을 목표로 한다. 이 장르의 핵심은 단순한 전투 승패가 아니라, 넓은 공간 안에서의 위치 선택, 이동 경로, 교전 회피 또는 개입 타이밍, 자원 확보, 그리고 변화하는 전장 조건에 대한 적응에 있다.

배틀로얄 장르에서 특히 중요한 구조는 시간이 지남에 따라 플레이 가능한 공간이 점차 축소된다는 점이다. 이러한 축소 구조는 플레이어에게 지속적인 이동 압박을 만들고, 경기 후반으로 갈수록 생존자들이 좁은 공간에 밀집하도록 만든다. 따라서 각 팀의 생존 가능성은 단순히 누적 kill 수나 damage량과 같은 사후 집계 통계만으로 충분히 설명되기 어렵다. 같은 전투 능력을 가진 팀이라도 현재 위치, 안전구역과의 거리, 주변 팀의 밀도, 이동해야 하는 거리, 직전 교전 상태, 다음 단계에서 요구되는 이동 부담에 따라 이후 생존 가능성은 크게 달라질 수 있다.

\textit{PUBG: Battlegrounds}는 이러한 시공간적 생존 구조가 명확하게 관측되는 대표적인 배틀로얄 게임이다. PUBG에서는 최대 100명의 플레이어가 전장에 진입해 무기와 장비를 획득하고, 시간이 지남에 따라 축소되는 공간 안에서 생존을 경쟁한다~\cite{pubgGameInfo}. 또한 PUBG의 match telemetry는 phase 변화, 플레이어 위치, 피해, knockdown, kill, elimination 등 경기 중 발생하는 다양한 사건을 시간 순서대로 기록한다~\cite{pubgTelemetryEvents}. 따라서 PUBG는 단순히 대중적인 게임 사례가 아니라, 이동 객체, 동적 공간 영역, 그리고 사건 발생이 결합된 시공간 telemetry 데이터를 분석하기에 적합한 연구 대상이다.

PUBG 경기는 시공간 데이터베이스 관점에서 세 가지 요소로 해석할 수 있다. 첫째, 각 팀은 시간에 따라 전장 위를 이동하는 객체이다. 둘째, 안전구역과 블루존은 시간에 따라 위치와 크기가 변화하는 동적 공간 영역이다. 셋째, 피해, knockdown, kill, elimination은 특정 시점과 위치에서 발생하는 사건이다. 본 연구는 이러한 원시 telemetry를 팀, phase, 이동, 공간 영역, 주변 팀 관계, 전투 사건을 포함하는 team-phase 수준의 시공간 event-history 데이터로 재구성한다.

기존의 PUBG 대상 공개 benchmark와 예측 연구는 주로 경기 종료 후 집계된 통계를 이용하여 final placement 또는 final rank를 예측하는 문제에 집중해 왔다~\cite{kagglePubgFinishPlacement,sen2021pubgFinalRank,lee2025pubgRank}. 이러한 접근은 전체 경기 결과를 요약적으로 예측한다는 점에서 유용하지만, 배틀로얄 경기의 단계적 진행 구조를 충분히 보존하지 못한다. 최종 순위 예측은 어떤 플레이어 또는 팀이 높은 순위를 얻었는지를 설명할 수는 있지만, 특정 팀이 언제부터 위험해졌는지, 어떤 phase에서 공간적 압박이 커졌는지, 안전구역 변화와 주변 팀 배치가 탈락 가능성을 어떻게 변화시켰는지를 직접적으로 분석하기 어렵다.

본 연구는 이 한계를 다루기 위해 \textit{phase-level next-phase elimination hazard prediction} 문제를 정의한다. 여기서 phase \(k\)는 경기 진행 중 관측되는 하나의 시간 구간이며, team-phase는 특정 팀이 특정 phase에서 관측되는 하나의 분석 단위를 의미한다. 본 연구는 phase \(k\)까지 관측된 정보, 즉 \(X_{i,\leq k}\)를 이용하여 팀 \(i\)가 다음 phase 구간에서 탈락할 조건부 위험을 예측한다.

구체적으로, phase \(k\)의 시작 시점을 \(a_k\), 종료 시점을 \(b_k\), 팀 \(i\)의 탈락 시점을 \(T_i\)라고 하자. 본 연구는 phase \(k\)의 종료 시점까지 생존한 팀에 대해, phase \(k\)까지 관측된 시공간 정보 \(X_{i,\leq k}\)를 사용하여 다음 phase 구간에서 탈락하는지를 예측한다. 결과 변수는 다음과 같이 정의된다.

\[
Y_{i,k}
=
\Ind\!\left\{b_k < T_i \leq b_{k+1}\right\}.
\]

이에 따라 phase-level next-phase elimination hazard는 다음과 같이 표현된다.

\[
h_{i,k}
=
P(b_k < T_i \leq b_{k+1}
\mid T_i > b_k, X_{i,\leq k}).
\]

이 정의는 동일 phase 내부에서 발생한 탈락을 같은 phase의 feature로 예측하는 것이 아니라, cut phase까지의 관측 정보를 바탕으로 다음 phase의 탈락 위험을 예측한다는 점에서 정보 시점과 예측 구간을 분리한다. 이는 이산 시간 event-history 및 survival analysis에서 각 시간 구간의 조건부 사건 발생 위험을 모델링하는 관점과 대응된다~\cite{allison1982discrete,singer1993discrete}. 즉, 본 연구의 관심은 전체 경기의 최종 순위를 한 번에 예측하는 것이 아니라, 경기 진행 중 각 phase에서 현재까지 생존한 팀이 다음 phase까지 생존하지 못할 조건부 위험을 반복적으로 추정하는 데 있다.

본 연구의 목적은 단순히 PUBG 경기의 승패나 최종 등수를 예측하는 데 있지 않다. 보다 넓게는 복잡한 telemetry 데이터를 시공간 event-history 데이터로 변환하고, 동적 공간 조건이 개체의 생존 또는 탈락 사건에 어떤 영향을 미치는지를 예측하고 해석하는 것이다. PUBG는 한 경기 안에 다수의 이동 객체가 존재하고, 안전구역이라는 동적 공간 제약이 명확하며, 탈락이라는 사건이 시간 순서대로 관측된다는 점에서 이러한 문제를 실증하기에 적합한 사례이다.

본 연구의 주요 기여는 다음과 같다.

\begin{itemize}
  \item 원시 PUBG match telemetry를 팀, phase, 이동, 공간 영역, 주변 팀 관계, 전투 사건을 포함하는 team-phase 수준의 시공간 event-history 데이터로 재구성한다.
  \item final placement 예측 대신 next-phase elimination hazard prediction 문제를 정의하여, 배틀로얄 경기의 단계적 생존 과정을 조건부 사건 예측 문제로 정식화한다.
  \item 안전구역 노출, 이동 부담, 주변 팀 밀도, 교전 상태와 같은 시공간 특징이 phase별 탈락 위험을 어떻게 설명하는지 분석한다.
  \item tabular feature 기반 모델과 graph-aware 모델을 비교하여, 주변 팀과의 거리, 밀도, 인접 관계와 같은 공간적 상호작용을 관계 구조로 표현하는 것이 phase-level hazard prediction에서 어떤 효과와 한계를 가지는지 검토한다.
\end{itemize}
